\section{Open AI System and Research Acceleration}
\label{sec:ai}
\subsection{Motivation}
A major barrier to audiology AI is the scarcity of harmonized, reusable clinical features. StressEar-AI addresses this by publishing a transparent feature schema, reproducible public-data code, and a clinician-verification protocol. This supports a broadly shared-benefit research model: investigators can reuse the pipeline, add local cohorts, and compare models without exposing identifiable patient data.

\subsection{Model Roles and Prohibitions}
The AI system has four restricted roles: (1) it harmonizes variable names and labels across public datasets; (2) it trains transparent baseline models for association and calibrated risk stratification; (3) in the clinical extension, an open medical language model (e.g., MedGemma) extracts structured fields from deidentified notes using a fixed JSON schema; and (4) it drafts explanations that cite only approved study documents or clinical guidelines. The system is explicitly prohibited from diagnosing an individual patient, assigning a definitive etiology, recommending or dosing steroids, or delaying urgent referral. These prohibitions are enforced in code and documented in a model card.

\subsection{Deterministic Red-Flag Safety Layer}
Before any probabilistic model output is surfaced, a deterministic rule layer screens for red flags derived from the SSNHL and tinnitus guidelines \citep{chandrasekhar2019ssnhl,tunkel2014cpg}: sudden or rapidly progressive hearing loss (especially unilateral), aural fullness with a discrete onset, single-sided tinnitus with asymmetry, vertigo, or focal neurologic signs. Any trigger overrides the model and returns an ``urgent audiology/otolaryngology referral'' message. The layer is evaluated primarily on recall (missed red flags are the costly error); its rules, versions, and test cases are unit-tested in the repository. This operationalizes the ``golden-window'' lesson from the motivating vignette: the framework must never let a probabilistic ``low-risk, likely stress'' estimate suppress timely referral.

\subsection{Schema-Constrained Extraction and Governance}
LLM extraction is constrained to a fixed schema with typed fields and controlled vocabularies; free-text generation is disallowed for structured outputs. Each extracted field is accompanied by a source span and a confidence flag, and every field is subject to clinician verification before entering analysis. Prompts, schema versions, model weights/version, temperature, and random seeds are logged. Extraction quality is reported as human-verified field accuracy, unsupported-output (hallucination) rate, and a clinician review of any potentially harmful output.

\subsection{Robustness and Trustworthiness}
Trustworthiness is pursued through eight design choices: (1) public-data baselines; (2) country-level external validation; (3) survey-weighted epidemiology; (4) interpretable models before complex models; (5) locked feature and outcome definitions; (6) clinician review of LLM-extracted variables; (7) subgroup calibration and error analysis; and (8) full run provenance (dataset versions, variable mappings, hyperparameters, prompts, seeds). Reproducibility is supported by pinned dependencies and deterministic seeds.

\subsection{Open-Science Contribution}
The repository includes executable analysis templates, a data dictionary, model cards, prompt templates, a red-flag test suite, and a completed reporting checklist. Because KNHANES files cannot be redistributed without official access procedures, the code specifies directory expectations and variable-mapping scripts rather than repackaging raw data. NHANES download helpers retrieve public files directly from CDC endpoints where available. A synthetic-data generator lets others exercise the full pipeline end-to-end without any restricted data, lowering the barrier to reuse and review.
